%% ma: L12-B08-0162
%% lop: 12
%% chuong: 2
%% bai: 8
%% dang: 12.8.52
%% loai: TN4
%% muc: VD
%% dap_an: "D"
%% nguon: "CĐ2.4.pdf giáo viên gửi 27/10/2026 (M07, Lê Minh Khôi – Nguyễn Anh Khôi, Chương II Chủ đề 4 Các bài toán về tâm tỉ cự), Câu 12"
%% kien_thuc: "Bài 8 (biểu thức tọa độ của tích vô hướng), lớp 10 (giá trị nhỏ nhất của tam thức bậc hai)"
%% trang_thai: cho-duyet
%% ly_do: "Dạng mới 12.8.52 đề xuất, chờ giáo viên duyệt"
\begin{ex}
Trong không gian $Oxyz$, cho ba điểm $A(1;-1;2)$, $B(-2;0;3)$, $C(0;1;-2)$. Gọi $M(a;b;c)$ là điểm thuộc mặt phẳng $(Oxy)$ sao cho biểu thức $S=\overrightarrow{MA}\cdot\overrightarrow{MB}+2\overrightarrow{MB}\cdot\overrightarrow{MC}+3\overrightarrow{MC}\cdot\overrightarrow{MA}$ đạt giá trị nhỏ nhất. Tính giá trị của biểu thức $T=12a+12b$.
\choice{$T=3$}{$T=-3$}{$T=1$}{\True $T=-1$}
\loigiai{$M(a;b;0)$: $S=6a^2+2a+6b^2-b-23$.

$S$ nhỏ nhất khi $a=-\dfrac16$, $b=\dfrac1{12}$. Vậy $T=-2+1=-1$.}
\end{ex}
